% =============================================================
% Capítulo 2. Marco teórico
% =============================================================

\chapter{Marco teórico}
\label{cap:marco}

La introducción ha planteado la construcción de mazos como un problema de optimización costoso y ha formulado la pregunta que vertebra el trabajo. Este capítulo establece los fundamentos necesarios para abordarla, siguiendo un hilo que parte del problema y desemboca en las herramientas para resolverlo. En primer lugar se describe Magic: The Gathering y se formaliza la construcción de mazos como problema de optimización combinatoria. A continuación se presentan los algoritmos genéticos, la metaheurística elegida para resolverlo, y el simulador Forge, que actúa como motor de evaluación. Por último se revisa el estado del arte en la construcción automática de mazos, que sitúa este trabajo respecto a la literatura previa e identifica la brecha que este trabajo aborda.

\section{Magic: The Gathering y la construcción de mazos}
\label{sec:mtg}

Magic: The Gathering es un juego de cartas en el que dos jugadores se enfrentan con sus respectivos mazos. Cada jugador comienza con una cantidad de puntos de vida y vence cuando reduce a cero los del rival, habitualmente mediante criaturas que atacan y hechizos que interactúan con el estado de la partida \cite{wotc_rules}. El recurso que permite jugar cartas es el maná, que se produce mayoritariamente con tierras y se divide en cinco colores: blanco, azul, negro, rojo y verde. Cada color se identifica de forma abreviada por su inicial en inglés: W (blanco), U (azul), B (negro), R (rojo) y G (verde). Los mazos se describen por los colores que emplean, y las combinaciones de dos colores reciben además nombres propios tomados del juego; por ejemplo, el par rojo-blanco (R/W) se denomina Boros, el azul-rojo (U/R) Izzet, el azul-blanco (U/W) Azorius y el azul-negro (U/B) Dimir. La etiqueta \emph{Mono-Red} designa un mazo de un solo color, el rojo.

La construcción del mazo está sujeta a reglas que el reglamento oficial fija con precisión \cite{wotc_rules}. En los formatos de \emph{constructed} (aquellos en los que cada jugador construye su mazo de antemano a partir de todas las cartas de las que dispone, por oposición a los formatos limitados, donde el mazo se arma con cartas abiertas en el momento), un mazo contiene un mínimo de sesenta cartas y no puede incluir más de cuatro copias de ninguna carta, con la única excepción de las tierras básicas, admitidas en cualquier número. Esta norma de las cuatro copias condiciona directamente el proceso de construcción: incluir una carta en un mazo equivale, casi siempre, a incluir cuatro copias de ella, ya que esa es la cantidad que maximiza la probabilidad de robarla cuando se necesita. La decisión del jugador no reside, por tanto, en cuántas copias incorporar, sino en si la carta entra o no en el mazo.

La base de tierras, o \emph{manabase}, constituye el subsistema que determina si el mazo puede pagar de forma consistente el maná de los colores que sus cartas exigen. Karsten formalizó este problema de manera cuantitativa y derivó heurísticas que relacionan el número de fuentes de un color con la demanda de dicho color en el mazo \cite{karsten_manabase}. La idea subyacente es que la probabilidad de disponer de un color en el momento oportuno crece con el número de fuentes que lo producen. El sistema de este trabajo incorpora una aproximación lineal a esas heurísticas: estima las fuentes necesarias de un color como una fracción de la demanda de símbolos de maná (\emph{pips}) de ese color, con un mínimo de tres fuentes,

\begin{equation}
\text{fuentes}(c) \;\approx\; \max\bigl(3,\; \lceil \kappa \cdot \text{pips}(c) \rceil \bigr), \qquad \kappa \approx 0.45,
\label{eq:karsten}
\end{equation}

donde $\text{pips}(c)$ es el número total de símbolos del color $c$ presentes en los costes de las cartas del mazo. Esta regla guía la reparación de la manabase de cada mazo evolucionado.

Los mazos de competición se agrupan en arquetipos según su estrategia. Los mazos de agresión, o \emph{aggro}, buscan una victoria temprana con criaturas baratas y una curva de maná baja. Los mazos de control son reactivos, juegan respuestas y cartas de coste elevado, y tratan de prolongar la partida hasta imponer su ventaja. En una posición intermedia, los mazos \emph{midrange} combinan criaturas eficientes con interacción. Esta clasificación, pese a su carácter simplificador, organiza el metajuego y orienta buena parte de las decisiones de construcción.

El formato Standard, sobre el que se centra este trabajo, restringe las cartas legales a las expansiones más recientes, publicadas aproximadamente en los últimos tres años, y renueva su contenido con cada rotación. Esa rotación obliga a reconstruir los mazos de forma periódica y mantiene un metajuego dinámico. La profundidad estratégica del juego admite una medida formal: se ha demostrado que el sistema de reglas de Magic es Turing-completo \cite{churchill2019}.

El tamaño del espacio de búsqueda puede acotarse para justificar la elección de un método aproximado. Un mazo de Standard se compone de alrededor de quince cartas no básicas distintas, cada una incluida como un conjunto de hasta cuatro copias, junto con una base de tierras. Considerando únicamente la elección de esas cartas distintas a partir de un catálogo de aproximadamente $4\,130$ cartas legales, el número de combinaciones posibles es del orden de

\begin{equation}
\binom{4130}{15} \approx 1.3 \times 10^{42}.
\label{eq:espacio}
\end{equation}

Esta cifra constituye una cota inferior, pues ignora tanto la variación en el número de copias como la composición de la manabase; el espacio completo es varios órdenes de magnitud mayor. Una magnitud de este orden descarta la enumeración exhaustiva como estrategia de resolución. La construcción de mazos se formula así como un problema de optimización combinatoria con restricciones duras y una función objetivo de evaluación costosa. Un problema con este perfil requiere una metaheurística de búsqueda~\cite{blum2003}; la elegida en este trabajo, el algoritmo genético, se presenta en la sección siguiente.

\section{Algoritmos genéticos}
\label{sec:ga}

Un algoritmo genético es una metaheurística de búsqueda inspirada en la evolución biológica. Holland lo formalizó en los años setenta como un modelo de adaptación en sistemas naturales y artificiales \cite{holland1975}, y Goldberg lo consolidó en la década siguiente como herramienta de búsqueda y optimización \cite{goldberg1989}. Su principio consiste en mantener una población de soluciones candidatas y hacerla evolucionar mediante operadores que emulan la selección natural, la reproducción y la mutación.

Formalmente, sea $X$ el espacio de soluciones candidatas y $f: X \rightarrow \mathbb{R}$ una función de fitness que asigna a cada individuo $x \in X$ un valor que mide su calidad. Cada individuo se codifica en una representación manipulable por el algoritmo, denominada genotipo. Partiendo de una población inicial $P_0 \subset X$, el algoritmo repite un ciclo generacional \cite{eiben2015}. La selección escoge los individuos que se reproducen, asignando mayor probabilidad a los de mayor fitness. El cruce combina el material de dos progenitores para formar descendientes. La mutación introduce variaciones aleatorias que preservan la diversidad de la población. Finalmente, un mecanismo de reemplazo determina qué individuos componen la generación siguiente. El proceso se repite hasta satisfacer un criterio de parada, normalmente un número máximo de generaciones o la consecución de un valor de fitness objetivo.

El equilibrio entre exploración y explotación gobierna el comportamiento de un algoritmo genético~\cite{crepinsek2013}. La exploración consiste en buscar en regiones nuevas del espacio de soluciones; la explotación, en refinar las soluciones de calidad ya encontradas. Una presión selectiva elevada combinada con una mutación baja favorece la explotación y aumenta el riesgo de convergencia prematura, situación en la que la población se estanca en un óptimo local antes de explorar lo suficiente. Una mutación elevada favorece la exploración, pero puede impedir que las buenas soluciones se consoliden. Diversas técnicas gestionan esta tensión. El elitismo preserva los mejores individuos de una generación a la siguiente para evitar su pérdida por azar \cite{mitchell1996}. La mutación adaptativa ajusta la intensidad de la variación en función del progreso del algoritmo~\cite{eiben1999}. La selección por torneo regula la presión selectiva mediante el tamaño de la muestra que compite por reproducirse~\cite{millergoldberg1995}.

Los algoritmos genéticos resultan apropiados para problemas con espacios de búsqueda amplios y discretos, sin gradiente disponible y con enumeración exhaustiva inviable, perfil que coincide con el caracterizado en la sección~\ref{sec:mtg}. No garantizan el óptimo global, pero suelen alcanzar soluciones de buena calidad en un tiempo razonable. El problema de este trabajo presenta, además, un agravante: la función de fitness no es una expresión cerrada, sino una simulación costosa y sujeta a ruido, lo que exige el motor de evaluación que se describe a continuación.

\section{Forge: motor de simulación}
\label{sec:forge}

La evaluación de la calidad de un mazo requiere jugarlo. Forge es una implementación de código abierto de las reglas de Magic que incorpora un oponente controlado por inteligencia artificial y permite disputar partidas completas sin intervención humana \cite{forge}. A diferencia de las heurísticas que puntúan un mazo a partir de su composición, Forge aplica el reglamento completo y resuelve cada partida jugándola, lo que convierte al simulador en un oráculo de fitness basado en el rendimiento real.

La elección de Forge no fue inmediata. En los primeros ensayos del proyecto se empleó Magarena, otro simulador de Magic de código abierto. Su catálogo de cartas, sin embargo, había quedado desactualizado respecto a las expansiones vigentes del formato Standard, lo que impedía representar el espacio de búsqueda de este trabajo. Forge, en cambio, mantiene un catálogo actualizado y ofrece una interfaz por línea de comandos adecuada para la evaluación automática de partidas, motivos por los que se adoptó como motor de simulación definitivo.

Forge ofrece un modo de simulación que enfrenta dos mazos y disputa una o varias partidas, con la inteligencia artificial gestionando ambos lados. Ese modo es el que la propuesta emplea para evaluar la población: cada mazo se enfrenta a otros y su tasa de victorias alimenta la función de fitness. La evaluación por simulación conlleva un coste elevado, dado que cada partida requiere varios minutos, y una componente de azar, ya que el reparto inicial de cartas y las decisiones de la inteligencia artificial introducen variabilidad. Presenta, asimismo, una limitación relevante para este trabajo: la inteligencia artificial de Forge ejecuta con desigual acierto las distintas estrategias y juega por debajo de su potencial los arquetipos de mayor complejidad táctica, un sesgo cuyas consecuencias se analizan más adelante.

Las cartas que componen el catálogo del sistema se obtienen de Scryfall, un servicio que ofrece una base de datos completa de Magic a través de una API pública \cite{scryfall}. De esta fuente se extraen los atributos de cada carta que el algoritmo necesita: coste de maná, tipo, colores, texto y características de combate.

Caracterizado el problema y presentadas la metaheurística y el motor de evaluación, resta situar el trabajo respecto a los enfoques previos. La sección siguiente revisa el estado del arte en la construcción automática de mazos.

\section{Estado del arte: construcción automática de mazos}
\label{sec:estado-arte}

La aplicación de la computación evolutiva a los juegos cuenta con una trayectoria amplia, en especial en la generación automática de contenido \cite{togelius2011}. Miles y Louis emplearon algoritmos genéticos para generar estrategias en juegos de estrategia en tiempo real \cite{miles2006}, mientras que Mahlmann, Togelius y Yannakakis los utilizaron sobre el juego de cartas Dominion para evolucionar conjuntos de cartas con fines de balanceo \cite{mahlmann2012}. Estos trabajos consolidan las metaheurísticas evolutivas como herramienta para los problemas de diseño y estrategia en juegos, si bien ninguno aborda la construcción de un mazo competitivo como objetivo directo.

La construcción automática de mazos posee identidad propia dentro de este campo, ya que combina un espacio de búsqueda combinatorio con una evaluación que idealmente requiere disputar partidas. El trabajo de referencia en juegos de cartas es el de García-Sánchez et al., que aplicó un algoritmo genético a la construcción de mazos en Hearthstone \cite{garciasanchez2016}. Su sistema representa cada mazo como una lista de cartas, emplea operadores de cruce y mutación a nivel de carta individual y evalúa el fitness mediante partidas simuladas contra un conjunto de mazos de referencia, demostrando que el método produce mazos competitivos frente a los diseñados por jugadores humanos. Hearthstone es, no obstante, un juego más simple que Magic: el maná se genera de forma automática y numerosas cartas están restringidas por clase, lo que reduce el espacio de búsqueda y elimina el problema de la manabase. Más allá de los enfoques evolutivos directos, el espacio de mazos se ha explorado con otras técnicas: estrategias evolutivas y algoritmos de calidad-diversidad para cartografiar el espacio de Hearthstone \cite{bhatt2018, fontaine2019, zhang2022}, la evolución del propio metajuego \cite{mesentier2019}, y el aprendizaje por refuerzo, tanto para la construcción por genes activos \cite{kowalski2020} como para la fase de \emph{draft} \cite{vieira2020}.

El antecedente más cercano a este trabajo es la tesis de máster de Bjørke y Fludal, que aborda la construcción de mazos en Magic mediante un algoritmo genético \cite{bjorke2017}. Su sistema evalúa los mazos disputando partidas con el motor Forge, el mismo simulador que emplea este Trabajo de Fin de Grado, y representa cada mazo como un vector de enteros sobre el que aplica cruce y mutación. La comparación con la presente propuesta revela tres diferencias relevantes. La primera afecta al formato: Bjørke y Fludal se restringen al formato sellado, con un repertorio de alrededor de noventa cartas, mientras que aquí se trabaja sobre el formato Standard completo, con miles de cartas. La segunda concierne a la granularidad de los operadores: su representación admite copias repetidas, pero sus operadores no preservan la estructura de conjuntos de cuatro copias que caracteriza la construcción de competición, que constituye precisamente la aportación central de este trabajo. La tercera atañe al esquema de evaluación: ellos puntúan cada mazo contra un conjunto fijo de seis mazos rivales, un esquema al que este trabajo también recurrió en una vía que finalmente se descartó.

Dos de los hallazgos de Bjørke y Fludal refuerzan los resultados de este trabajo. En primer lugar, observaron que sus poblaciones convergían con rapidez a una monocultura debido al reducido tamaño poblacional, un fenómeno análogo a la convergencia cromática que también aparece en los experimentos de este trabajo. En segundo lugar, y de manera más significativa, detectaron que el motor Forge favorece a unos arquetipos sobre otros: en uno de sus experimentos, el mazo construido por un jugador experto rindió por debajo del mazo generado por el algoritmo, lo que les llevó a concluir que Forge no reproduce con fidelidad el juego de un experto humano. Ese sesgo del simulador coincide con el que este trabajo encontró al explorar esa misma vía; su aparición en un estudio independiente confirma que se trata de una limitación del motor y no de un artefacto de esta implementación.

En el entorno de la propia Escola Politècnica Superior de Gandia existe un precedente directo: el Trabajo de Fin de Grado de Mira Abad, que aplicó neuroevolución a Magic con un objetivo distinto \cite{mira2021}. Su trabajo evoluciona la estrategia de juego de un agente, modelada como un conjunto de máquinas de estados, empleando el motor Magarena para simular las partidas. La diferencia con la presente propuesta reside en el objeto de la evolución: Mira Abad evoluciona el modo en que se juega un mazo, mientras que aquí se evoluciona qué mazo se construye. Ambos enfoques son complementarios y comparten el uso de la simulación como función de evaluación. Más allá de las técnicas evolutivas, el propio juego de Magic se ha abordado con búsqueda en árbol: Cowling, Ward y Powley aplicaron Monte Carlo Tree Search con determinización por conjuntos para gestionar la información oculta y la aleatoriedad inherentes al juego \cite{cowling2012}, si bien su objetivo es decidir jugadas y no construir el mazo. En una línea diferente, se han empleado redes neuronales para analizar las propias cartas de Magic a partir de su texto e ilustración \cite{zilio2018}, y en Hearthstone se ha mejorado el nivel de juego de los agentes combinando MCTS con aprendizaje supervisado \cite{swiechowski2018}.

La tabla~\ref{tab:comparativa} sintetiza los trabajos revisados y los contrasta con la presente propuesta.

\begin{table}[H]
\centering
\caption{Síntesis comparativa de los trabajos previos y su relación con la propuesta.}
\label{tab:comparativa}
\small
\begin{tabular}{p{2.3cm}p{1.7cm}p{2.2cm}p{3.0cm}p{3.0cm}}
\toprule
\textbf{Referencia} & \textbf{Dominio} & \textbf{Técnica} & \textbf{Limitación principal} & \textbf{Diferencia con la propuesta} \\
\midrule
Miles \& Louis (2006) & RTS & Algoritmo genético & No aborda construcción de mazos & Problema y dominio distintos \\
Mahlmann et al. (2012) & Dominion & Algoritmo genético & Objetivo de balanceo, no de competición & Genera mazos competitivos \\
García-Sánchez et al. (2016) & Hearthstone & Algoritmo genético & Juego simplificado; sin manabase & Standard completo; operadores a nivel de pack \\
Bjørke \& Fludal (2017) & Magic (sellado) & Algoritmo genético + Forge & Repertorio reducido; operadores no pack-aware & Standard; norma de 4 copias; torneo Swiss \\
Mira Abad (2021) & Magic & Neuroevolución & Evoluciona la estrategia, no el mazo & Evoluciona la construcción del mazo \\
\bottomrule
\end{tabular}\\[2pt]
{\footnotesize Fuente: elaboración propia.}
\end{table}

Un aspecto transversal a esta literatura es el tratamiento de la evaluación estocástica. Cuando el fitness se obtiene de partidas simuladas, su valor presenta varianza, y los algoritmos evolutivos disponen de estrategias específicas para operar bajo esa incertidumbre \cite{jinbranke2005}. Este trabajo afronta dicha incertidumbre mediante la repetición de partidas y un esquema de emparejamiento eficiente.

En síntesis, la revisión confirma dos conclusiones. La construcción evolutiva de mazos es un enfoque validado, pero poco explorado en Magic y casi siempre sobre formatos o juegos de espacio reducido. La evaluación por simulación, aunque ofrece la medida más fiel del rendimiento, arrastra el sesgo del motor que la ejecuta. Sobre esta base, la propuesta de este trabajo se diferencia de los antecedentes en tres elementos combinados: el uso del formato Standard completo, unos operadores genéticos que respetan la norma de las cuatro copias y un sistema de torneo que mantiene acotado el coste de la evaluación. Antes de detallar esa propuesta, el capítulo siguiente expone la metodología con la que se desarrolló el trabajo.
